% Chapter 5: Proposed Methodology
% Needs in the preamble:
%   \usepackage{amsmath,booktabs,tikz}
%   \usetikzlibrary{positioning,arrows.meta,fit,backgrounds}
%   \usepackage{algorithm}
%   \usepackage{algpseudocode}

\chapter{Proposed Methodology}

\section{Design (architecture) or conceptual framework}

The proposed system is a QoE-aware, self-adaptive educational assistant for NCERT Class~10 Science. It combines
three ideas: retrieval-augmented generation (RAG) grounded in a lesson video, a hybrid edge--cloud execution path in
which each question can be answered by a small local model or a large cloud model, and a closed adaptation loop in
which every learning session is measured and stored as an \emph{experience} that guides the runtime configuration of
later sessions. A proactive, concept-level coach sits on top of this loop and turns the learner's test results into
learning goals and practice.

Figure~\ref{fig:architecture} shows the architecture. It is organised in six layers.

\begin{enumerate}
  \item \textbf{Presentation layer.} A Streamlit web interface for learners (tests, video choice, questions, practice
        checks, mastery panel) and a FastAPI backend. Optional speech input and output make the interface accessible.
  \item \textbf{Orchestration layer.} One LangGraph state graph runs a learning session as a sequence of nodes. The
        graph pauses at every point that needs the learner (an \emph{interrupt}) and resumes when the learner replies,
        so a session survives page reloads and its full state is inspectable.
  \item \textbf{Knowledge layer.} Recommended YouTube videos are the learning resources. The chosen video's transcript
        is cleaned, chunked and embedded into a FAISS index; corrective RAG (CRAG) retrieves and grades evidence from
        it for every question.
  \item \textbf{Model layer.} A local small language model (gemma3:1b on Ollama, the \emph{edge}) and a cloud large
        language model (gpt-oss-120b on Groq, the \emph{cloud}) behind a router. Embeddings are computed locally
        (nomic-embed-text).
  \item \textbf{Adaptation layer.} Context monitoring (network round-trip time, CPU and memory load, learner
        competency, previous QoE), session evaluation (QoE, SRE, utility, reward) and policy selection over an
        experience store indexed with FAISS.
  \item \textbf{Coaching layer.} A concept vocabulary per topic, a concept tracker that estimates mastery per learner
        and concept, session planning, a rule-based coach, a pool of formative questions and a parallel post-test.
\end{enumerate}

All persistent data (learner profiles, sessions, assessments, question banks, concept mastery, experiences) is kept in
a relational database (SQLite), and every node is traced with Langfuse for analysis.

\begin{figure}[htbp]
\centering
\begin{tikzpicture}[
  font=\small,
  box/.style={draw, rounded corners=2pt, align=center, minimum height=0.9cm, text width=3.0cm},
  wide/.style={box, text width=6.6cm},
  layer/.style={draw, dashed, rounded corners=4pt, inner sep=6pt},
  arr/.style={-{Latex[length=2mm]}, thick},
  node distance=0.55cm and 0.6cm]

  % presentation
  \node[box] (learner) {Learner\\(text or speech)};
  \node[box, right=of learner] (ui) {Streamlit UI\\FastAPI backend};

  % orchestration
  \node[wide, below=0.9cm of $(learner)!0.5!(ui)$] (graph)
       {LangGraph session agent\\pre-test $\rightarrow$ plan $\rightarrow$ video $\rightarrow$ questions
        $\rightarrow$ coach $\rightarrow$ post-test $\rightarrow$ evaluation};

  % coaching and knowledge
  \node[box, below left=0.9cm and -1.4cm of graph] (coach)
       {Coaching layer\\concept tracker,\\session plan, coach,\\formative pool};
  \node[box, below right=0.9cm and -1.4cm of graph] (knowledge)
       {Knowledge layer\\YouTube recommender,\\transcript $\rightarrow$ FAISS,\\CRAG retrieval};

  % model layer
  \node[box, below=0.9cm of knowledge] (router) {LLM router\\(policy route)};
  \node[box, below left=0.7cm and 0.2cm of router, text width=2.4cm] (edge) {Edge: local SLM\\gemma3:1b};
  \node[box, below right=0.7cm and 0.2cm of router, text width=2.4cm] (cloud) {Cloud: LLM\\gpt-oss-120b};

  % adaptation
  \node[box, below=0.9cm of coach] (monitor) {Context monitor\\RTT, CPU, memory,\\competency, last QoE};
  \node[box, below=0.6cm of monitor] (eval) {Session evaluation\\QoE, SRE, utility,\\reward};
  \node[box, below=0.6cm of eval] (store) {Experience store\\(DB + FAISS)};
  \node[box, left=0.8cm of eval] (select) {Policy selection\\kNN + $\varepsilon$-greedy\\(P1--P5)};

  \node[box, below=1.2cm of $(edge)!0.5!(cloud)$, text width=6.0cm] (db)
       {Relational DB (SQLite) and Langfuse traces};

  \begin{scope}[on background layer]
    \node[layer, fit=(monitor)(eval)(store)(select), label={[font=\footnotesize\itshape]above left:Adaptation loop}] {};
    \node[layer, fit=(router)(edge)(cloud), label={[font=\footnotesize\itshape]above:Hybrid edge--cloud execution}] {};
  \end{scope}

  \draw[arr] (learner) -- (ui);
  \draw[arr] (ui) -- (graph);
  \draw[arr] (graph) -- (coach);
  \draw[arr] (graph) -- (knowledge);
  \draw[arr] (knowledge) -- node[right, font=\footnotesize]{evidence} (router);
  \draw[arr] (router) -- (edge);
  \draw[arr] (router) -- (cloud);
  \draw[arr] (graph.west) -| node[pos=0.75, left, font=\footnotesize]{context} (select.north);
  \draw[arr] (monitor) -- (eval);
  \draw[arr] (eval) -- (store);
  \draw[arr] (store.west) -| (select.south);
  \draw[arr] (select) to[bend right=10] node[below, font=\footnotesize]{policy $a^\ast$} (router.south west);
  \draw[arr, dashed] (edge.south) -- (db);
  \draw[arr, dashed] (cloud.south) -- (db);
\end{tikzpicture}
\caption{Architecture of the QoE-aware self-adaptive hybrid edge--cloud educational assistant.}
\label{fig:architecture}
\end{figure}

\paragraph{Key concepts.}
\begin{itemize}
  \item \textbf{Runtime policy.} A predefined configuration $a = (\text{route}, k, r)$: where answers are generated
        (local, cloud, or local with escalation), how many chunks are retrieved ($k$), and how many corrective
        retrieval retries CRAG may make ($r$). Five policies are defined (Table~\ref{tab:policies}).
  \item \textbf{Context (state).} The measurable situation in which a session runs:
        $s = (\mathit{rtt}, \mathit{cpu}, \mathit{mem}, \mathit{competency}, \mathit{QoE}_{\mathit{last}})$.
  \item \textbf{Experience.} A completed transition $(s, a, R, s')$ of one session: the context, the policy executed,
        the reward and the measured outcome.
  \item \textbf{Concept.} A unit of the chapter (for example \emph{dominant and recessive traits}); every test item,
        practice question and learner question is tagged with one concept, and mastery is tracked per concept.
\end{itemize}

\paragraph{Assumptions.}
(i) The learning resource is a single lesson video at a time, and answers should be grounded in its transcript.
(ii) The local model is cheaper and private but weaker; the cloud model is stronger but depends on the network.
(iii) Conditions that matter for the choice of policy (network delay, device load, learner level) can be measured at
the start of a session. (iv) A session is the unit of adaptation: one policy is applied per session and its effect is
measured at the end of the session.

\section{Proposed approach and procedure}

\subsection{Overview of a learning session}

One learning session follows the procedure below. Each step is a node (or a small group of nodes) of the session
graph.

\begin{enumerate}
  \item \textbf{Pre-test.} The learner chooses a topic and answers five multiple-choice questions drawn from a reviewed
        15-question bank, one question per concept.
  \item \textbf{Session plan.} The concepts of the wrongly answered questions, and earlier weak concepts, become two to
        four learning goals.
  \item \textbf{Video recommendation.} YouTube videos are searched with a query built from the topic and the goals,
        ranked, and shown to the learner, who chooses one.
  \item \textbf{Knowledge preparation.} The chosen video's transcript is fetched, chunked and indexed.
  \item \textbf{Context measurement and policy selection.} The context $s$ is measured and a policy $a^\ast$ is
        selected from similar past experiences.
  \item \textbf{Question answering.} For each learner question, CRAG retrieves evidence and the router generates the
        answer according to $a^\ast$.
  \item \textbf{Coaching.} After each answer the question is tagged with a concept, mastery is updated, and the coach
        suggests one next step (for example a short practice check).
  \item \textbf{Post-test and feedback.} The learner answers five post-test questions parallel to the pre-test ones and
        rates the session.
  \item \textbf{Evaluation and storage.} QoE, SRE, utility and reward are computed, and the experience is stored for
        future sessions.
\end{enumerate}

The following subsections describe each stage.

\subsection{Learning content}

For each chapter a \emph{content package} is prepared from the NCERT textbook: 8--10 concepts; 15 pre-test questions;
15 post-test questions, each \emph{parallel} to one pre-test question (same concept and difficulty, different
question); 3--4 formative (practice) questions per concept, kept separate from both tests; and 3--5 candidate
videos. Every question has four options, the answer, an explanation, one concept, a difficulty and the textbook
section and page that support it. Packages are checked automatically before loading: grounding of every item in the
cited page, absence of long passages copied from the textbook, balanced answer positions, distinct concepts among the
first five pre-test questions, and an independent ``blind'' check in which a second model answers and tags every
question without seeing the key. Disagreements are not accepted or rejected automatically; they are flagged for human
review. Five chapters are prepared: Life Processes, Control and Coordination, How do Organisms Reproduce?,
Heredity, and Our Environment.

\subsection{Pre-test and session planning}

The pre-test shows the first five bank questions, which cover five different concepts. Each answer is recorded as an
attempt on its concept, and mastery is estimated with the Laplace (rule-of-succession) estimate
\begin{equation}
  m_{c} = \frac{n^{\text{correct}}_{c} + 1}{n^{\text{attempts}}_{c} + 2},
  \label{eq:mastery}
\end{equation}
which starts at $0.5$ for an unseen concept and moves towards the observed accuracy as evidence accumulates. The
session plan selects two to four goals in priority order: concepts missed in this pre-test, then concepts with low
mastery from earlier sessions, then, only to reach the minimum, concepts never practised.

\subsection{Resource recommendation}

The recommendation query combines the topic, the goal concepts and a hint for the learner's level (for example
``explained simply'' for beginners). Candidates from the YouTube Data API are filtered by hard constraints (privacy,
minimum duration, previously completed videos) and ranked in two stages. A cheap first stage scores titles and
descriptions by semantic similarity to the learning objective; the second stage fetches transcripts only for the
most promising candidates and scores how well their content covers the objective. Three recommendation strategies
trade speed for evidence (metadata only, or transcripts for the top four or eight candidates). Because the session
answers questions from the transcript, only videos with a transcript in the learner's language are recommended.

\subsection{Knowledge preparation and corrective retrieval}

The chosen video's transcript is cleaned, split into overlapping chunks of 120 words (25-word overlap) and embedded
into a per-video FAISS index; the index is reused when the video is chosen again. For a question $q$, CRAG retrieves
the $k$ most similar chunks, grades them, and computes a retrieval score. If the score is below the acceptance
threshold $\tau = 0.65$ and retries remain, the query is rewritten (filler words removed, topic terms added) and
retrieval is repeated; the attempts are merged. Retrieval stops when the score reaches $\tau$ or the policy's retry
budget $r$ is spent.

\subsection{Hybrid edge--cloud generation}

The router generates the answer from the retrieved chunks according to the policy's route
(Table~\ref{tab:policies}). P1--P4 are fixed configurations and serve as baselines. P5 is a confidence-aware cascade:
\begin{itemize}
  \item if the final retrieval score lies in the borderline band $[0.55, 0.65)$, the evidence exists but is weakly
        matched, so the cloud model answers directly;
  \item otherwise the local model answers, and a fast grounding check (a small cloud model) measures the share of the
        answer's claims supported by the retrieved chunks; if it is below $0.75$, the cloud model answers again and its
        answer replaces the local draft.
\end{itemize}
Below the band (score $< 0.55$) the question is outside the video's content, and the local model states that the
video does not cover it, at no cloud cost.

\begin{table}[htbp]
\centering
\caption{Predefined runtime policies (the actions of the adaptation loop).}
\label{tab:policies}
\begin{tabular}{llccl}
\toprule
Policy & Route & $k$ & $r$ & Behaviour \\
\midrule
P1 local-light (default) & local & 3 & 0 & local model only \\
P2 local-thorough & local & 5 & 1 & local model only \\
P3 cloud-light & cloud & 3 & 0 & cloud model only \\
P4 cloud-thorough & cloud & 5 & 1 & cloud model only \\
P5 local-escalate & escalate & 5 & 1 & local first, cloud when needed \\
\bottomrule
\end{tabular}
\end{table}

\subsection{Proactive coaching}

After every answer the learner's question is matched to a concept: by name or alias when exactly one concept is
named, otherwise by embedding similarity with a confidence threshold. The concept is marked as practised, and the
coach chooses exactly one next step by the first matching rule:
\begin{enumerate}
  \item \emph{formative check}: the question concerned a goal concept with mastery below $0.5$;
  \item \emph{another video}: a goal concept is not covered by the current video;
  \item \emph{review}: the last answer was escalated, failed the grounding check, or was weakly retrieved;
  \item \emph{post-test}: every goal has mastery of at least $0.75$;
  \item \emph{follow-up}: the next goal not yet asked about, otherwise the weakest goal.
\end{enumerate}
A suggestion ignored twice in a row is muted. Formative questions come from the concept's practice pool and are never
test questions, so practice does not leak the post-test. Rules rather than a learned model are used deliberately at
this stage: they are transparent to the learner, need no training data, and give a clear baseline for comparing a
learned coach later.

\subsection{Post-test}

The post-test asks, for each pre-test question, its parallel question. Because the two forms test the same concept
at the same difficulty with different wording, a higher post-test score reflects learning of the concept rather than
memory of the question. Post-test questions and their answers are stored separately from the pre-test bank and are
never included in pre-test data. The normalized learning gain is
\begin{equation}
  g = \frac{\text{post} - \text{pre}}{1 - \text{pre}}.
  \label{eq:gain}
\end{equation}

\subsection{Session evaluation}

All measurements are normalised to $[0,1]$ with higher meaning better (costs such as latency are inverted). Quality
of experience combines four groups:
\begin{align}
  \mathit{EQ} &= \alpha_1 \mathit{LG} + \alpha_2 \mathit{PT} + \alpha_3 \mathit{RA} + \alpha_4 \mathit{LS}, \\
  \mathit{AQ} &= \beta_1 \mathit{RQ} + \beta_2 \mathit{LQ} + \beta_3 \mathit{SQ} + \beta_4 \mathit{TQ}, \\
  \mathit{NQ} &= \gamma_1 \mathit{PL} + \gamma_2 \mathit{NR} + \gamma_3 \mathit{RU} + \gamma_4 \mathit{AV}, \\
  \mathit{UX} &= \delta_1 \mathit{SR} + \delta_2 \mathit{FS} + \delta_3 \mathit{CR}, \\
  \mathit{QoE} &= w_{\mathit{EQ}}\,\mathit{EQ} + w_{\mathit{AQ}}\,\mathit{AQ} + w_{\mathit{NQ}}\,\mathit{NQ}
                  + w_{\mathit{UX}}\,\mathit{UX},
\end{align}
where the educational group uses learning gain (LG), post-test score (PT), recommendation acceptance (RA) and
resource satisfaction (LS); the answer-quality group uses retrieval quality (RQ), response quality (LQ) and speech
quality (SQ, TQ); the network group uses answer latency (PL), network reliability (NR), device utilisation (RU) and
availability of a model route (AV); and the user-experience group uses the session rating (SR), feedback sentiment
(FS) and completion (CR). System resource efficiency measures how much of a per-session budget was used:
\begin{equation}
  \mathit{SRE} = \eta_1 \Bigl(1-\tfrac{u_{\mathit{cpu}}}{u^{\max}_{\mathit{cpu}}}\Bigr)
               + \eta_2 \Bigl(1-\tfrac{u_{\mathit{mem}}}{u^{\max}_{\mathit{mem}}}\Bigr)
               + \eta_3 \Bigl(1-\tfrac{E}{E^{\max}}\Bigr)
               + \eta_4 \Bigl(1-\tfrac{B}{B^{\max}}\Bigr),
\end{equation}
with each term clipped to $[0,1]$, where $E$ is an energy estimate and $B$ the bytes transferred. SRE alone would
reward doing less, so it is always combined with QoE. The utility of an observed state and the reward of a
transition $(s, a, s')$ are
\begin{align}
  U &= \lambda_1\, \mathit{QoE} + \lambda_2\, \mathit{SRE}, \\
  R &= \mu_1 \bigl(\mathit{QoE}' - \mathit{QoE}\bigr) + \mu_2 \bigl(\mathit{SRE}' - \mathit{SRE}\bigr),
\end{align}
with $\lambda_1 = \mu_1 = 0.7$ and $\lambda_2 = \mu_2 = 0.3$. Every weight group is non-negative and sums to one.
When a measurement is unavailable, the weights of the available ones in its group are renormalised and the share of
weight backed by data (\emph{coverage}) is reported, so a score never claims more evidence than it has. As an
alternative to expert weights, the Entropy Weight Method can derive weights from the variation in stored data.

\subsection{Experience-based policy selection}

The context is embedded as a vector $\mathbf{x}(s) \in [0,1]^5$ (round-trip time divided by a maximum and clipped;
CPU and memory load; competency; previous QoE). Completed transitions are stored in the database and their context
vectors in an exact FAISS index. For a new session, the $K = 5$ nearest past contexts are retrieved. For each policy
$a$, only neighbours that executed $a$ contribute, weighted by a Gaussian kernel of their distance $d_i$:
\begin{equation}
  \hat{U}(s, a) = \frac{\sum_{i \in N_a(s)} w_i\, U_i}{\sum_{i \in N_a(s)} w_i},
  \qquad w_i = \exp\!\Bigl(-\frac{d_i^2}{2\sigma^2}\Bigr), \quad \sigma = 0.25,
  \label{eq:estimate}
\end{equation}
where $U_i$ is the utility measured after neighbour $i$. A policy has an estimate only when the total weight of its
neighbours is at least $0.1$, so evidence from very different contexts is not trusted. The selection is
$\varepsilon$-greedy (Algorithm~\ref{alg:policy}). This is case-based reasoning rather than reinforcement learning: no
model is trained; past outcomes are looked up. The predicted utility of the chosen policy is stored alongside the
utility observed at the end of the session, which allows the accuracy of the predictions to be measured.

\begin{algorithm}[htbp]
\caption{Experience-based policy selection}
\label{alg:policy}
\begin{algorithmic}[1]
\Require context $s$, experience store $\mathcal{D}$, policies $\mathcal{A}$, default policy $a_0$ (P1),
         $K$, $\sigma$, exploration rate $\varepsilon$, minimum experiences $n_{\min}$
\If{$|\mathcal{D}| < n_{\min}$} \State \Return $a_0$ \Comment{cold start}
\EndIf
\State $N \gets$ the $K$ nearest neighbours of $\mathbf{x}(s)$ in $\mathcal{D}$
\For{each $a \in \mathcal{A}$}
  \State $N_a \gets \{\, i \in N : a_i = a \,\}$
  \If{$\sum_{i \in N_a} w_i \ge 0.1$} \State compute $\hat{U}(s,a)$ with Eq.~\eqref{eq:estimate}
  \EndIf
\EndFor
\If{no policy has an estimate} \State \Return $a_0$
\EndIf
\State $a^{\ast} \gets \arg\max_a \hat{U}(s,a)$ \Comment{ties go to $a_0$}
\If{$\mathrm{rand}() < \varepsilon$} \State $a^{\ast} \gets$ a random policy other than $a^{\ast}$ \Comment{exploration}
\EndIf
\State store $\hat{U}(s, a^{\ast})$ as the prediction for this session
\State \Return $a^{\ast}$
\end{algorithmic}
\end{algorithm}

\subsection{Implementation}

The system is implemented in Python~3.11. The session graph uses LangGraph with a checkpointer, so each pause for the
learner is a resumable state. The backend uses FastAPI, the interface Streamlit, and persistence SQLite through
SQLAlchemy with versioned, additive migrations. Retrieval uses FAISS with locally computed embeddings
(nomic-embed-text on Ollama); the local language model is gemma3:1b on Ollama and the cloud model is gpt-oss-120b on
Groq. Videos come from the YouTube Data API, transcripts from youtube-transcript-api. Speech output uses the Kokoro
text-to-speech model; speech input is optional. Every node and model call is traced in Langfuse. The design choices
were made for transparency and reproducibility: fixed, inspectable policies; deterministic metrics computed from
stored measurements; case-based selection whose every decision can be explained by its neighbours; and rule-based
coaching. The implementation is covered by automated tests, and the main configuration values are listed in
Table~\ref{tab:config}.

\begin{table}[htbp]
\centering
\caption{Main configuration values.}
\label{tab:config}
\begin{tabular}{lr}
\toprule
Parameter & Value \\
\midrule
Pre-test / post-test questions per session & 5 / 5 \\
Learning goals per session & 2--4 \\
Transcript chunk size / overlap (words) & 120 / 25 \\
CRAG acceptance threshold $\tau$ & 0.65 \\
P5 borderline retrieval band & $[0.55, 0.65)$ \\
P5 grounding-check threshold & 0.75 \\
Coach: low mastery / mastered & 0.5 / 0.75 \\
Neighbours $K$, kernel width $\sigma$ & 5, 0.25 \\
Exploration rate $\varepsilon$, minimum experiences & 0.2, 3 \\
Utility and reward weights (QoE, SRE) & 0.7, 0.3 \\
\bottomrule
\end{tabular}
\end{table}
